% ---------------------------------------------------------------------
% 2.12 Modelado del lenguaje y predicción de palabras
% ---------------------------------------------------------------------
\section[MODELADO DEL LENGUAJE Y PREDICCIÓN DE PALABRAS]{Modelado del lenguaje y predicción de
palabras}
\label{sec:lenguaje}

\subsection{Predicción de palabras en la comunicación asistida}

Aun con una detección perfecta de la P300, escribir carácter por carácter en un P300 Speller es un
proceso lento, pues cada selección requiere varios segundos. La predicción de palabras, ampliamente
utilizada en los sistemas de comunicación aumentativa, reduce el número de selecciones porque le
ofrece al usuario una lista de palabras probables a partir de lo que ya escribió, y con una sola
selección más puede completar la palabra. Ryan et al. integraron
este mecanismo en un P300 Speller agregando a la matriz una región con las palabras sugeridas, y
encontraron que los participantes completaban las frases en menos tiempo que con el deletreador
convencional, aunque con una exactitud por selección ligeramente menor \cite{ryan2011}. El
beneficio de estos sistemas suele medirse mediante el ahorro de selecciones,
\begin{equation}
\text{AS} = \frac{S_{\text{sin}} - S_{\text{con}}}{S_{\text{sin}}}\times100\%,
\label{eq:ahorro}
\end{equation}
donde $S_{\text{sin}}$ y $S_{\text{con}}$ son el número de selecciones necesarias para escribir un
mismo texto sin y con predicción, respectivamente.

\subsection{Modelos de lenguaje}

Un modelo de lenguaje asigna una probabilidad a cada secuencia de palabras, y con ella se puede
estimar cuál es la palabra, o el carácter, más probable dado lo que se ha escrito. Los modelos más simples
son los de $n$-gramas, que suponen que cada palabra depende únicamente de las $n-1$ palabras
anteriores. Sus probabilidades se estiman contando cuántas veces aparece cada secuencia en un
corpus de texto \cite{jurafsky2025}. Debido a que están limitados a un contexto corto y fijo,
han sido reemplazados progresivamente por modelos neuronales, entre ellos los modelos de lenguaje
de gran escala (LLM), entrenados con enormes cantidades de texto \cite{brown2020}.

\subsection{Integración del lenguaje en el P300 Speller}

La información del lenguaje puede incorporarse al P300 Speller de dos maneras: como sugerencias
explícitas que el usuario selecciona, o como una probabilidad previa que se combina con los
puntajes del clasificador para decidir el carácter. Speier et al. siguieron la segunda vía,
integrando un modelo de lenguaje en el proceso de clasificación dinámica, y reportaron mejoras
tanto en la exactitud como en la tasa de transferencia de información \cite{speier2012}. Más
recientemente, Hong et al. propusieron ChatBCI, un P300 Speller que utiliza un LLM para predecir
palabras según el contexto de la frase, logrando una reducción del tiempo promedio de escritura de
62.14\% \cite{hong2026}.

